\scanpage{118}
% L1–2
\uline{N\textsuperscript{o} 1} Préschémas \unsure{cpts|cpcts} et loc. cpts. Anneaux flasques.

\uline{N\textsuperscript{o} 2} \deleted{Systèmes projectifs}\inserted{Proschémas} flasques \deleted{\unsure{de préschémas et} \unsure{constructibisation} \illegible{3}}

% L3–4
\uline{N\textsuperscript{o} 3} \deleted{Espaces \deleted{\unsure{ind}} schématiques Ind = \unsure{Espaces|ensembles}}\inserted{espaces schématiques compacts} Espaces schématiques. [Catégorie \unsure{semblable pour|véritable pour|engendrée par}] les rec. ouverts \ldots{} ]
\margin{Morphismes quasi-compacts $\rightarrow$ \unsure{séparation|opérations} propres}
% slanted

% L5
\uline{N\textsuperscript{o} 4} Opérations sommes, \deleted{\unsure{sous}}\inserted{\unsure{cas des quasi-compacts|des proschémas} $\simeq$ sommes de compacts} et $\underleftarrow{\lim}$ \deleted{\unsure{finies}} des espaces schématiques

% L6
\uline{N\textsuperscript{o} 5} \deleted{\unsure{Esp. quasi-schématiques}} Espaces quasi-schématiques \ldots\ldots

% L7–9: cancelled diagram
\deleted{\diagram{\resizebox{\linewidth}{!}{\begin{tikzcd}[ampersand replacement=\&,column sep=large]
\left(\substack{\text{Préschémas quasi-compacts}\\\text{quasi-séparés}}\right)\arrow[r,"\alpha"]
\&\left(\substack{\text{Esp. schématiques}\\\text{compacts}}\right)
 \arrow[r,"{\substack{\inserted{\text{équiv.}}\\\deleted{\unsure{\text{fidèle}}}\\\approx}}"]
\&\left(\operatorname{Pro}(\mathrm{Sch})\ \substack{\text{flasques}\\\deleted{\unsure{\text{finis}}}\ \text{q.cpts}}\right)
 \arrow[r,"{\substack{\text{fid, }\inserted{\text{submersif}}\\\deleted{\unsure{\text{ess. surjectif}}}}}"]
\&\left(\substack{\text{schémas compacts}\\\text{réduits}}\right)
 \arrow[l,bend left=24,"{\text{pl. fid. }\unsure{\sigma|\gamma}}"']
 \arrow[lll,bend left=32,"{\text{pl. fid.}}"]
\end{tikzcd}}}}

% L10
\uline{N\textsuperscript{o} 6} \unsure{pts|pbs|Pb} épais. \unsure{Espaces|Espace} de Néron. \deleted{\illegible{1}}

% L11–15
\diagram{\resizebox{\linewidth}{!}{\begin{tikzcd}[ampersand replacement=\&,column sep=huge]
\substack{(\text{Préschémas quasi-compacts}\\\text{quasi-séparés})\\[3pt] X}
 \arrow[r,"{\alpha}","{\substack{\text{submersif}\\\text{[fidèle sur les}\\\text{réduits].}}}"']
 \arrow[rr,bend left=24,"{\sigma}"]
\&\substack{\text{(Cons cpt)}\\ \shortparallel\\(\operatorname{Pro}(\mathrm{Sch})\ \text{flasques})\\[3pt]\mathscr Y}
 \arrow[r,"{\beta=\underleftarrow{\lim}}","{\substack{\text{submersif}\\\text{fidèle}}}"']
\&\substack{(\text{schémas compacts}\\\text{réduits})\\[3pt]K}
 \arrow[l,bend left=24,"{\beta'\ \text{pl. fid}}"']
 \arrow[ll,bend left=33,"{\sigma'\ \text{pl. fid. }\unsure{\text{imm. can.}|\text{imm. ouv.}|\text{inv. can.}|\text{immersion}}}"']
\end{tikzcd}}}
\margin{$\nearrow$ [pleinement fidèle également sur les ``de type fini'' \unsure{en un sens|au sens}\ \unsure{précisé\ldots|fonc\ldots}]}
% L15–16
$\beta'(K)={}$ système proj trivial défini par $K$

% L17–18
\[
\begin{aligned}
\operatorname{Hom}(K,\sigma X)&\xrightarrow{\sim}\operatorname{Hom}(\sigma'K,X)
&\operatorname{Hom}(K,\beta(\alpha X))&=\operatorname{Hom}(\beta'K,\alpha X)\\
\operatorname{Hom}(K,\beta\mathscr Y)&\xrightarrow{\sim}\operatorname{Hom}(\beta'K,\mathscr Y)
\end{aligned}
\]
% archive sheet number: 117
